\documentclass[10pt,a4paper]{article}
\usepackage[T1]{fontenc}
\usepackage{lmodern}
\usepackage[margin=23mm]{geometry}
\usepackage{amsmath,amssymb,booktabs,tabularx,microtype}
\usepackage[hidelinks,hypertexnames=false]{hyperref}
\pdfinfoomitdate=1
\pdftrailerid{}
\pdfsuppressptexinfo=15
\setlength{\parindent}{0pt}
\setlength{\parskip}{5pt}
\newcommand{\kb}{k_{\mathrm B}}
\newcommand{\mh}{m_{\mathrm H}}
\begin{document}
\begin{center}
{\Large A connected hydrogen escape atmosphere}\\[4pt]
Equations, code paths and conditions for scientific use
\end{center}

\section{A connected calculation is not an equivalence claim}
The native entry point is
\path{escape.atmosphere.HydrogenAtmosphereModel.spectrum}.
It solves internal radiation, hydrogen populations, temperature and a
transonic wind, then passes those same populations to the observer.
No MCEG or Cloudy import enters this calculation.
Independent Cloudy/sunbather experiments live under \path{tools/}.
A successful solve does not establish equivalence to Zhang et al.~\cite{zhang}:
chemical composition, resonance redistribution, the dense lower atmosphere
and observation processing need separate comparisons. A pure-H energy equation
cannot be assumed adequate for a solar-composition atmosphere.

\section{Mass, pressure and momentum}
For the native pure-H closure, with $x=n_p/n_{\rm H}$,
\begin{equation}
 n_e=n_p=x n_{\rm H},\quad \mu=\frac{1}{1+x},\quad
 \rho=\mh n_{\rm H}=\frac{\dot M}{4\pi r^2v},\quad
 a^2=\frac{P}{\rho}=\frac{\kb T}{\mu\mh}.
 \label{eq:mass}
\end{equation}
The wind satisfies
\begin{equation}
 \left(v-\frac{a^2}{v}\right)\frac{dv}{dr}
 =\frac{2a^2}{r}-\frac{da^2}{dr}-\frac{GM_p}{r^2}+\kappa r,
 \qquad \kappa=\frac{3GM_\star}{a_{\rm orb}^3}.
 \label{eq:wind}
\end{equation}
The regular critical point has $v=a$ and a vanishing right-hand side.
\path{atmosphere_wind.wind_from_profiles} supplies $T/\mu$ to the existing
implicitly differentiated Parker solver with its analytic isothermal
initialization. Thus the spatial derivative of composition is retained.
A cubic Hermite curve interpolates $\log(T/\mu)$ against $\log r$,
with zero endpoint slopes and constant extension to the wind boundary
($32R_p$). This is a boundary assumption, not a solved outer temperature.
Its sensitivity and grid convergence require their own checks.

\section{Radiation and statistical equilibrium}
The atom retains resolved $1s,2s,2p,3s,3p,3d$ and statistical shells through
the configured $n_{\max}$. At each radius,
\begin{equation}
 0=\sum_j(n_jR_{ji}-n_iR_{ij})+n_en_p\alpha_i+n_e^2n_p\beta_i
   -n_i\Gamma_i-n_en_iq_{i,\mathrm{ion}},
 \qquad \sum_i n_i+n_p=n_{\rm H}.
 \label{eq:se}
\end{equation}
Here $R_{ij}$ includes radiative and collisional transitions. Populations
and radiation depend on each other:
\begin{equation}
 \mu_\Omega\frac{dI_\nu}{dr}=j_\nu-\chi_\nu I_\nu,
 \quad J_\nu=\frac12\int_{-1}^1I_\nu\,d\mu_\Omega,
 \quad \overline J_{ul}=\int\phi_{ul}(\nu)J_\nu\,d\nu.
 \label{eq:rt}
\end{equation}
The native formal solver uses a static, plane-parallel slab and complete
redistribution (CRD). It integrates homogeneous cells exactly and uses their
cell-average intensity for absorption. For thickness $\Delta s$ and
$t=\chi_\nu\Delta s$, let $w_1=(1-e^{-t})/t$ and $w_2=(1-w_1)/t$:
\begin{equation}
 I_{\rm out}=I_{\rm in}e^{-t}+j_\nu\Delta s\,w_1,
 \qquad \langle I\rangle=I_{\rm in}w_1+j_\nu\Delta s\,w_2.
 \label{eq:cell}
\end{equation}
Series limits retain transparent emission and derivatives. Cell absorption
minus emission agrees with the face-flux difference. Negative extinction is
outside this formal solver's current domain. Accurate profile normalization
alone does not ensure that escape through far damping wings is resolved;
angular, spatial and frequency refinement remain separate requirements.

The complete atmosphere uses a normally incident stellar beam. Its boundary
array is the equivalent intensity $I_{\rm eq}=2J_\star$, giving normal flux
$F_{\rm beam}=4\pi J_\star$. For stellar surface flux and dilution,
$J_{\star,\nu}=W F^\star_\nu/\pi$. An isotropic hemisphere with the same
incident mean has half that normal flux. The lower incoming continuum is
independently specified; it is not automatically identical to Cloudy's
\texttt{double optical depths} boundary.

Cloudy C23 uses a different line closure: a special Ly$\alpha$
escape/destruction treatment, a core escape approximation for other Lyman
lines, and wings for subordinate lines. CRD transport is not automatically
equivalent. Two-photon emission currently escapes as a cooling channel;
its internal reabsorption is another explicit adequacy question.

\section{Thermal advection and expansion}
The baseline follows sunbather's local statistical-equilibrium thermal
closure, without population advection~\cite{sunbather}:
\begin{align}
 q_{\rm adv}&=-\rho v\frac{d}{dr}
      \left(\frac{\kb T}{(\gamma-1)\mu\mh}\right),
 &q_{\rm exp}&=\frac{\kb T v}{\mu\mh}\frac{d\rho}{dr},\notag\\
 q_{\rm adv}+q_{\rm exp}+\Gamma_{\rm th}-\Lambda_{\rm th}&=0,
 &\gamma&=5/3.
 \label{eq:thermal}
\end{align}
All terms are erg\,cm$^{-3}$\,s$^{-1}$. Decreasing density gives
$q_{\rm exp}<0$. Local radiative equilibrium alone omits both flow terms.
Advected chemical binding energy belongs to a different reactive-flow
closure and is not added here.

\path{hydrogen_continuum.thermal_energy_exchange} evaluates kinetic heating
without subtracting enormous trapped radiative cycles. Photoionization
supplies photon energy above threshold; recombination removes kinetic energy;
collisions exchange kinetic and atomic internal energy. For bound-bound
collisions,
\begin{equation}
 q_{\rm coll}=n_e\sum_{ij}n_iq_{ij}(E_i-E_j).
 \label{eq:collheat}
\end{equation}
At SE this kinetic exchange equals the full photon energy transfer. Away
from SE the difference is atomic internal-energy production,
$-\sum_i\chi_i\dot n_i$. This identity is independently tested.
An extra two-photon thermal loss would double count its energy pathway.
Thermal tolerances use actual kinetic-process turnover, not the much larger
trapped-photon luminosity.

The optional low-temperature law is
$\Upsilon(T)=\Upsilon(T_0)T/T_0$ below the existing
$T_0=0.5\,\mathrm{eV}/\kb$ node, following Cloudy's extrapolation.
It extends only to the 4000 K ionization-fit boundary and is not new atomic
data. The join is continuous with a derivative kink; older rate APIs still
reject extrapolation by default.

An optional static number of dense bottom cells imposes supplied temperature
and Saha/Boltzmann populations. This LTE photospheric patch does not solve
those cells' thermal balance. It requires a demonstrated thermalized join
and pressure/domain checks. A cold solution is never silently clipped into
the atomic domain, and the join does not move during differentiation.

\section{Root, derivatives and observer}
The root uses $\log(T/T_{\rm init})$ and independent $\log(n_i/n_p)$;
a softmax imposes nucleus and charge conservation. Every residual evaluation
recomputes $\mu$, wind, density, radiation, rates and thermal-flow terms.
Radiation continuation at fixed initial temperature precedes release of
the thermal equation. For a converged state $\mathcal F(y,\theta)=0$,
\begin{equation}
 \frac{dy}{d\theta}=-\mathcal F_y^{-1}\mathcal F_\theta.
 \label{eq:ad}
\end{equation}
JAX differentiates this equation and the wind root, not Newton's iteration
count. Failed roots return flags and NaN profiles, never a fallback spectrum.

\path{nlte_forward.HydrogenNLTEForwardModel.spectrum_from_profiles} consumes
the same populations, sums the three actual H$\alpha$ branches and reuses
ExoJAX intensity transfer, disk integration and instrumental convolution.
Continuum emissivity converts from per Hz to per wavenumber with a factor
$c$. For the complete atmosphere, H-minus is a construction-time option
shared by internal radiation/energy and the observer. Its existing trace-LTE
and temperature limits remain; this is not easychem chemistry.

The separate \path{observation.process_phase_spectra} operator takes explicit
exposure phases, orbital shifts, night labels, masks and weights. It shifts
planet flux to the stellar frame, adds the stellar reference, normalizes the
continuum, divides by the nightly master, masks contaminated pixels, returns
to the planet frame and coadds. The master must include all its exposures,
including those excluded from the final phase bin. The supplied stellar
reference and planet flux must use the same continuum units. No weights are
inferred from the existence of the operator; interpolation across masked
holes is excluded. Exposure-time and wavelength-bin integration are not
implemented. The archived observational example uses a declared reconstruction,
not the authors' unpublished exact processing recipe.

\section{Required evidence}
\path{mceg/uses/yapeng_atmosphere.yaml} selects scoped Assessments for the
actual spectrum entry point. Wiring and derivatives are distinct from
matched Cloudy references, composition, lower-atmosphere, resolution,
redistribution and observation-processing adequacy. A FollowUp records work,
not scientific support. The current records and retained numerical results
determine readiness; this document's existence does not.

A controlled Cloudy composition comparison makes one missing contribution
concrete. At $r=1.1R_p$ in the fixed-temperature, fixed-density PHOENIX
experiment, the net radiative thermal exchange is approximately
$-4.68\times10^{-4}$ erg\,cm$^{-3}$\,s$^{-1}$ for H+He, including H-minus
chemistry, but $+3.02\times10^{-2}$ for the specified solar mixture. The sign
change rejects composition-independent thermal adequacy for that experiment.
In the corresponding atomic-only solar diagnostic, Fe II, Ca II, Mg II,
Si II and C I account for about 89\% of local heating. These are pumping,
cascade and collisional de-excitation pathways, not only photoionization.
This is evidence against a hydrogen-only thermal reduction of the solar
reference; the paper does not specify its exact abundances. Agreement of
hydrogen source functions at a prescribed temperature would not repair the
missing energy budget. The source receipts, physical scope and unresolved
numerical controls are retained in the MCEG review.

\begin{thebibliography}{9}
\bibitem{zhang} Zhang et al. (2026), Sections IV--V,
\href{https://arxiv.org/abs/2606.23899v1}{arXiv:2606.23899v1}.
\bibitem{sunbather} Linssen et al. (2024), Sections 2.1--2.2.3,
\href{https://arxiv.org/abs/2404.12775}{arXiv:2404.12775}.
\bibitem{cloudy} Cloudy C23.01 source:
\path{hydrocollid.cpp}, \path{iso_level.cpp}, \path{rt_line_one.cpp},
\path{rt_ots.cpp}; \url{https://gitlab.nublado.org/cloudy/cloudy}.
\end{thebibliography}
\end{document}
